\chapter{Introducción}

En este trabajo práctico se aborda la programación de microcontroladores con núcleo ARM 
Cortex-M3, específicamente el \textbf{STM32F103C8T6} presente en la placa \textbf{BluePill}, 
utilizando un enfoque de nivel bajo o \textbf{bare-metal}. El objetivo es comprender el 
funcionamiento del hardware y las herramientas de desarrollo sin capas de abstracción 
intermedias que oculten los detalles de implementación.

\section{Hardware Utilizado}

Para el desarrollo de las actividades se emplearon los siguientes elementos:

\begin{itemize}
    \item \textbf{Placa BluePill}: Basada en el microcontrolador STM32F103C8T6. Cuenta con 
    una frecuencia de reloj de $72 \MHz$, $64 \kB$ de memoria Flash y $20 \kB$ de RAM.
    \item \textbf{Programador ST-Link V2}: Interfaz que permite la grabación y depuración 
    del microcontrolador a través del protocolo SWD.
\end{itemize}

\section{Software y Herramientas}

Se utilizó el ecosistema de herramientas estándar de GNU para sistemas embebidos:

\begin{itemize}
    \item \textbf{GNU ARM Toolchain}: Colección de compiladores, enlazadores y utilidades 
    para la arquitectura ARM (\texttt{arm-none-eabi-gcc}).
    \item \textbf{Make}: Herramienta para la automatización del proceso de compilación.
    \item \textbf{OpenOCD}: Software para la comunicación con el programador y el 
    microcontrolador.
    \item \textbf{GDB}: Depurador estándar para inspeccionar el estado del programa en 
    tiempo de ejecución.
\end{itemize}
